\documentclass[11pt]{article}

% SciPost-like concluded proof template for NIMA reports.
% Replace article with the official SciPost class if it is available locally.
\usepackage[a4paper,margin=1in]{geometry}
\usepackage{amsmath,amssymb,amsthm,mathtools}
\usepackage{hyperref}
\usepackage{enumitem}
\usepackage{booktabs}

\newtheorem{theorem}{Theorem}
\newtheorem{lemma}[theorem]{Lemma}
\newtheorem{proposition}[theorem]{Proposition}
\newtheorem{corollary}[theorem]{Corollary}
\newtheorem{definition}[theorem]{Definition}
\newtheorem{remark}[theorem]{Remark}

\title{Title of the Mathematical Research Report}
\author{Author names}
\date{\today}

\begin{document}
\maketitle

\begin{abstract}
State exactly what is proved, falsified, or left conditional. Name the main
claim id, proof attempt ids, formal theorem names, and whether trusted imports
or numerical certificates are used.
\end{abstract}

\section{Introduction}
Explain the mathematical question, why the strengthened or impossible statement
matters, and what changes relative to the starting conjecture or literature.

\section{Related Work and Imported Theorems}
Separate literature theorems, manuscript imports, numerical certificates, and
new results. Every trusted import should appear in the trust ledger in
Section~\ref{sec:trust-ledger}.

\section{Main Results}
State all systems, dimensions, maps, states, constants, and quantifiers.

\begin{theorem}[Main theorem or counterexample]
Precise statement.
\end{theorem}

\section{Proof Overview}
Give the proof architecture before technical details. Match the outline to the
NIMA DAG: claims, attempts, bridge gaps, and failed routes.

\section{Proof}
Give the complete proof step by step. Keep named lemmas aligned with NIMA
claim ids or theorem-card ids when possible.

\section{Experiments and Falsified Routes}
Document numerical searches, failed sublemmas, seeds, dimensions, scripts,
artifacts, and certificates. Explain exactly what each negative result rules
out.

\section{Formalization and Trust Ledger}
\label{sec:trust-ledger}

Inspect the exact NIMA session export and any Verify Lean receipts before filling this table.
Generated scientific results are uncertified; record the actual scope of formal verification separately.

\begin{center}
\begin{tabular}{ll}
\toprule
Item & Value \\
\midrule
NIMA claim id & \texttt{clm\_...} \\
Proof attempts & \texttt{att\_...} \\
Lean theorem names & \texttt{Namespace.theoremName} \\
Lean axiom class counts & \texttt{\{...\}} \\
Trusted literature imports & \texttt{LiteratureImports....} \\
Trusted manuscript imports & \texttt{ManuscriptGapAxioms....} \\
Numerical certificates & \texttt{ncert\_...} \\
Open bridge gaps & \texttt{fgap\_...} \\
Known failed routes & \texttt{att\_...}, \texttt{obs\_...} \\
\bottomrule
\end{tabular}
\end{center}

\section{Errors or Corrections in Source Papers}
List any source-paper errors, ambiguous constants, missing hypotheses, OCR
issues, or notation conflicts identified during the proof.

\section{Conclusion}
State what is now possible, what is impossible, and what remains conditional or
open.

\appendix

\section{Notation, Shorthands, Neologisms, and Symbols}
This appendix is mandatory. Define every project-local shorthand and every term
whose meaning is not standard in the cited literature. Include witness names,
bridge names, imported theorem aliases, entropy shorthands, norm conventions,
and all symbols used in the main theorem.

\begin{description}[leftmargin=2.5cm,style=nextline]
\item[Example shorthand] Precise definition and where it is first used.
\item[Example symbol] Domain, codomain, hypotheses, and convention.
\end{description}

\section{NIMA Packet Extract}
Include the relevant packet summary, source spans, theorem cards, formalization
gaps, and dependency edges.

\section{Lean Axiom Audit}
Include the raw or summarized \texttt{\#print axioms} output and its
recorded verification status.

\section{Additional Numerical Certificates}
Include reproducibility commands, seeds, dimensions, hashes, and interval
bounds.

\end{document}
